\documentclass[12pt]{exam}
\usepackage[margin=0.1in]{geometry}

\usepackage{amsmath,amssymb}
\usepackage{graphicx}
\usepackage{enumitem}
\usepackage{multicol}
\usepackage{xcolor}
\usepackage{tikz}
\usetikzlibrary{arrows.meta,positioning,shapes.callouts}

\newcommand{\tfq}[3]{%
  \question[#1]\textbf{T/F:} #2%
  % Only show the choices when printing answers/keys
  \ifprintanswers
    \par\begin{oneparchoices}
      #3
    \end{oneparchoices}
  \fi
}



\makeatletter
\newcommand{\TFQ@choices}{} % temp holder

\newenvironment{tfqverb}[2]{%
  \def\TFQ@choices{#2}% store choices for use at end
  \question[#1]\textbf{T/F:}%
}{%
  \ifprintanswers
    \par\begin{oneparchoices}\TFQ@choices\end{oneparchoices}
  \fi
}
\makeatother

%\printanswers
% optional styling for the correct choice:
\CorrectChoiceEmphasis{\bfseries}    % make the correct choice bold
% or e.g. \CorrectChoiceEmphasis{\color{red}\bfseries}

\usepackage{fancyvrb}
\newsavebox{\verbbox}

\begin{document}


\begin{center}
\textbf{\Large Review 2 — Database Management Systems}

\textbf{CSC 3300 \quad Spring 2026}
\end{center}

\begin{questions}
%---------------------------------TRUE/FALSE---------------------------------
\section*{Part I — True/False}
\tfq{2}{$\{course\_id, title\}$ is a candidate key for the relation \texttt{course}.}{\choice True \CorrectChoice False}
\tfq{3}{It is possible for the University Database to have exactly one tuple in each of its tables.}{\CorrectChoice True \choice False}
\tfq{3}{If the attribute \textit{grade} was part of the current primary key of the relation \texttt{takes}, then a student could receive more than one grade for a given course section in a given semester and year.}{\CorrectChoice True \choice False}
\tfq{3}{A tuple is inserted into the \texttt{instructor} relation only if the value of the \textit{ID} attribute of this tuple is equal to the value of the \textit{i\_id} attribute of some tuple in the \texttt{advisor} relation.}{\choice True \CorrectChoice False}
\begin{tfqverb}{3}{\choice True \CorrectChoice False}
Below query can return a tuple with value 0 for the attribute \textit{classroom\_no}.

\begin{verbatim}
select building, count(distinct room_number) as classroom_no
from classroom
group by building;
\end{verbatim}
\end{tfqverb}

  \tfq{3}{A query with an aggregate function in the 'select' part of it and without a 'group by' statement returns only one tuple.}{\CorrectChoice True \choice False}
  \begin{tfqverb}{3}{\CorrectChoice True \choice False}
  Below query returns courses \((\mathit{course\_id}, \mathit{title}, \mathit{dept\_name}, \mathit{credits})\) that have more than one direct prerequisite.
  \begin{verbatim}
    select * 
    from course 
    where (select count(course_id) 
           from prereq 
           where prereq.course_id = course.course_id) >= 2;
  \end{verbatim}
  \end{tfqverb}
  \begin{tfqverb}{3}{\CorrectChoice True \choice False}
  In the result of the provided query, there might be instructors whose information (\textit{ID}, \textit{name}, \textit{dept\_name}, \textit{salary}) is not included.
  \begin{verbatim}
  select * from instructor natural left outer join teaches;
  \end{verbatim}
  \end{tfqverb}

\begin{tfqverb}{3}{\choice True \CorrectChoice False}
The DBMS will utilize an index to retrieve the result of the following query:\par
\begin{verbatim}
select * 
from course
where course_id = 3300
  \end{verbatim}
HINT: MYSQL: explain\par\vspace{0.5em}
  \end{tfqverb}
%---------------------------------MC---------------------------------
\section*{Part II — Multiple choice (choose one)}

\question[3] Assume that there is only one tuple in the \textbf{section} relation and no tuples in both \textbf{takes} and \textbf{student} relations. What is the minimal number of tuples that have to be added to the database so that the query below returns 2 tuples?

\begin{verbatim}
select * from student natural join takes;
\end{verbatim}

\begin{choices}
  \choice 2
  \choice 3
  \CorrectChoice 4
  \choice 5
\end{choices}


\question[3]
What does the following SQL query return?\par\vspace{0.5em}
\begin{verbatim}
select i_id, count(s_id) 
from advisor 
group by i_id
\end{verbatim}
\begin{choices} 
    \choice For each instructor the number of students he advises.
    \choice For each advisor the number of students he advises.
    \choice For each student the number of his advisors.
    \CorrectChoice None of the other options.
\end{choices}

\question[4]
Which query is not equivalent to the remaining ones?\par\vspace{0.5em}
\begin{choices} 
    \choice \begin{verbatim}
select d.dept_name, d.building, c.course_id, c.title
from department as d, course as c
where d.dept_name = c.dept_name
    \end{verbatim}
    \choice \begin{verbatim}
select d.dept_name, d.building, c.course_id, c.title
from department as d join course as c using (dept_name)
    \end{verbatim}
    \choice \begin{verbatim}
select d.dept_name, d.building, c.course_id, c.title
from department as d inner join course as c on (d.dept_name=c.dept_name)
    \end{verbatim}
    \CorrectChoice \begin{verbatim}
select d.dept_name, d.building, c.course_id, c.title
from department as d right outer join course as c on (d.dept_name=c.dept_name)
    \end{verbatim}
\end{choices}

\question[4]
Choose the SQL query that is guaranteed to encounter an error due to issues with its structure or syntax.\par\vspace{0.5em}
\begin{choices} 
    \choice \begin{verbatim}
select course_id, sec_id, count(distinct ID)
from takes
group by course_id, sec_id
    \end{verbatim}
    \choice \begin{verbatim}
select count(distinct ID)
from takes
group by course_id, sec_id
    \end{verbatim}
    \choice \begin{verbatim}
select course_id, sec_id
from takes
group by course_id, sec_id
    \end{verbatim}
    \CorrectChoice \begin{verbatim}
select course_id, sec_id, semester, year, count(distinct ID)
from takes
group by course_id, sec_id
    \end{verbatim}
\end{choices}


%---------------------------------MS---------------------------------
\section*{Part III — Multiple select (select all that apply)}
\question[4] Choose superkeys of the relation \texttt{section}.
\begin{choices}
  \choice $\{\textit{building}, \textit{room\_no}\}$
  \CorrectChoice $\{\textit{course\_id}, \textit{sec\_id}, \textit{semester}, \textit{year}\}$
  \CorrectChoice $\{\textit{course\_id}, \textit{sec\_id}, \textit{semester}, \textit{year}, \textit{building}, \textit{room\_no}\}$
  \CorrectChoice $\{\textit{course\_id}, \textit{sec\_id}, \textit{semester}, \textit{year}, \textit{building}, \textit{room\_no}, \textit{time\_slot\_id}\}$
\end{choices}

\question[4] Consider the following specification of user requirements for a Library Database: For each book copy, the following data is stored: title, edition, number of pages, and barcode. When a member checks out a book copy, its barcode is scanned and linked to the member’s account. When the book copy is returned, the barcode is scanned again, and the book is marked as available for borrowing.

Consider storing title (\textit{title}), edition (\textit{edition}), number of pages (\textit{pages\_no}), and barcode (\textit{barcode}) for a book copy in a single table.

Which of the following statements are true? (Select all that apply.)

\begin{choices}
  \CorrectChoice $\{\textit{barcode}, \textit{title}\}$ is a super key.
  \choice $\{\textit{barcode}, \textit{title}\}$ is a candidate key.
  \CorrectChoice $\{\textit{barcode}\}$ is a candidate key.
  \CorrectChoice If multiple copies of the same book exist, the title may be stored multiple times, leading to data duplication.
\end{choices}

\question[6] Students (their \(\mathit{ID}, \mathit{name}, \mathit{dept\_name}, \mathit{tot\_cred}\)) who don't have advisors.
\begin{choices}

  \choice
\begin{verbatim}
with temp(ID) as
  (select s_id
   from advisor
   where i_id is null)
select *
from student natural join temp;
\end{verbatim}

  \choice
\begin{verbatim}
select ID, name, dept_name, tot_cred
from student natural join advisor
where i_id is null;
\end{verbatim}

  \CorrectChoice
\begin{verbatim}
select ID, name, dept_name, tot_cred
from student left outer join advisor
  on (student.ID = advisor.s_id)
where advisor.s_id is null;
\end{verbatim}

  \choice
\begin{verbatim}
select *
from student
where ID not in (select i_id from advisor);
\end{verbatim}

  \choice
\begin{verbatim}
select ID, name, dept_name, tot_cred
from student natural join advisor natural join instructor
where ID is null;
\end{verbatim}

\end{choices}

\question[4] Retrieve all course sections (identified by \((\mathit{course\_id}, \mathit{sec\_id}, \mathit{semester}, \mathit{year})\)) where the total enrollment exceeds the seating capacity of the assigned classroom.

\begin{choices}

  \CorrectChoice
\begin{verbatim}
select s.course_id, s.sec_id, s.semester, s.year
from section s
join takes t
  on s.course_id = t.course_id
 and s.sec_id    = t.sec_id
 and s.semester  = t.semester
 and s.year      = t.year
join classroom c
  on s.building = c.building
 and s.room_number = c.room_number
group by s.course_id, s.sec_id, s.semester, s.year, c.capacity
having count(t.ID) > c.capacity;
\end{verbatim}

  \CorrectChoice
\begin{verbatim}
with enroll(course_id, sec_id, semester, year, students_no) as (
  select course_id, sec_id, semester, year, count(*)
  from takes
  group by course_id, sec_id, semester, year
)
select s.course_id, s.sec_id, s.semester, s.year
from enroll e
join section s using (course_id, sec_id, semester, year)
join classroom c
  on s.building = c.building and s.room_number = c.room_number
where e.students_no > c.capacity;
\end{verbatim}

  \choice
\begin{verbatim}
with temp(course_id, sec_id, semester, year, students_no) as (
  select course_id, sec_id, semester, year, count(ID)
  from section natural join takes
  group by course_id, sec_id, semester, year
)
select course_id, sec_id, semester, year
from temp natural join classroom
where capacity < students_no;
\end{verbatim}

  \choice
\begin{verbatim}
select course_id, sec_id, semester, year
from section natural join takes
group by course_id, sec_id, semester, year
having count(*) > capacity;
\end{verbatim}

\end{choices}

\end{questions}
\end{document}